
% Estilo dos slides de conteudo

\definecolor{bordaSlide}{HTML}{8EAECB}
\definecolor{pontoInativo}{HTML}{C4D5E4}
\definecolor{pontoAtivo}{HTML}{104B80}

\setbeamercolor{normal text}{fg=black,bg=white}
\setbeamerfont{normal text}{
  size=\fontsize{10}{13}\selectfont
}

\setbeamercolor{frametitle}{fg=azulPrincipal}
\setbeamerfont{frametitle}{
  size=\fontsize{14}{17}\selectfont,
  series=\mdseries
}

\setbeamersize{
  text margin left=0.94cm,
  text margin right=0.94cm
}

% Titulo superior

\setbeamertemplate{frametitle}{%
  \nointerlineskip
  \vspace*{1.02cm}%
  \begin{beamercolorbox}[
    wd=\paperwidth,
    leftskip=0.73cm,
    rightskip=0.73cm
  ]{frametitle}%
    \usebeamerfont{frametitle}%
    \insertframetitle
  \end{beamercolorbox}%
  \vspace*{0.08cm}%
}

% Marcadores dos topicos

\setbeamertemplate{itemize item}{
  \raisebox{0.07ex}{
    \textcolor{azulPrincipal}{
      \scriptsize$\blacktriangleright$
    }
  }
}

\setbeamertemplate{itemize subitem}{
  \textcolor{azulPrincipal}{
    \scriptsize$\triangleright$
  }
}

\setlength{\leftmargini}{1.2em}

% Moldura e paginacao

\newcommand{\DesenharMolduraSlide}{%

  \ifnum\value{framenumber}>0\relax

    \begin{tikzpicture}[
      remember picture,
      overlay
    ]

      \coordinate (SW) at
        (current page.south west);

      \coordinate (NE) at
        (current page.north east);

      % Moldura principal

      \draw[
        draw=bordaSlide,
        line width=.70pt,
        rounded corners=.36cm
      ]
        ($(SW)+(.36cm,.32cm)$)
        rectangle
        ($(NE)+(-.36cm,-.32cm)$);

      % Quantidade automatica de bolinhas

      \pgfmathtruncatemacro{\TotalPontos}{
        max(1,\inserttotalframenumber)
      }

      \pgfmathsetmacro{\EspacoPontos}{
        min(.275,5.8/\TotalPontos)
      }

      % Bolinhas do progresso

      \foreach \i in {1,...,\TotalPontos} {

        \pgfmathsetmacro{\Deslocamento}{
          (\i-(\TotalPontos+1)/2)*\EspacoPontos
        }

        \ifnum\i=\value{framenumber}

          \fill[pontoAtivo]
            ($(current page.south)+
              (\Deslocamento cm,.64cm)$)
            circle [radius=.061cm];

        \else

          \fill[pontoInativo]
            ($(current page.south)+
              (\Deslocamento cm,.64cm)$)
            circle [radius=.061cm];

        \fi
      }

      % Espaco branco para a numeracao

      \fill[white]
        ($(current page.south)+
          (-.37cm,.22cm)$)
        rectangle
        ($(current page.south)+
          (.37cm,.42cm)$);

      % Numero da pagina

      \node[
        font=\fontsize{8.5}{10}\selectfont,
        text=azulPrincipal
      ]
        at ($(current page.south)+
          (0,.32cm)$)
        {
          \insertframenumber/
          \inserttotalframenumber
        };

    \end{tikzpicture}%

  \fi
}

\setbeamertemplate{background}{
  \DesenharMolduraSlide
}
